% Source: 01 Properties_of_Matter 26(1).docx; formatting in layout.tex.
\HOne{Further Applications of the Particle Model}

In Chapter 1.4, the particle model was used to explain the states of matter and several physical properties. The same model can also help us understand temperature, the spontaneous mixing of particles, the movement of water through membranes, and the structure and separation of mixtures.


\HTwo{Temperature}

\textbf{Temperature} is related to the motion of particles and therefore to the average kinetic energy of particles.
At higher temperatures, particles have more kinetic energy on
average, although individual particles may move at different speeds.

\[E_{\mathrm{kin}}=\tfrac12 \cdot m \cdot v^2\]

\textbf{Absolute temperature} is measured in kelvin (K). The Kelvin scale begins at \textbf{absolute zero} (0 K = -273.15 $^{\circ}$C), the theoretical lowest possible temperature. For a gas, the average kinetic energy of the particles increases in direct proportion to the absolute temperature.


% Maxwell-Boltzmann-Verteilung: zentral definiert in figures.tex
\TempChart

As temperature increases, the distribution shifts towards higher
speeds and becomes broader. Slow and fast particles are present
at every temperature, but a larger proportion of particles move
quickly at higher temperatures.


\begin{minipage}[c]{0.13\linewidth}
  \includegraphics[width=\linewidth]{images/image34.png}
\end{minipage}
\begin{minipage}[c]{0.83\linewidth}
  \href{https://phet.colorado.edu/sims/html/gas-properties/latest/gas-properties_all.html?locale=en}
  {\textcolor{DeepBlue}{\underline{Here}}}
  , an online simulation of particles at different
  temperatures can be found.
\end{minipage}\hfill



\HTwo{Diffusion and Osmosis}

Particles in liquids and gases move continuously and randomly.
If their concentration differs between two regions, this motion
produces a net movement from higher to lower concentration.
This process is called \textbf{diffusion}. It continues until
the concentration is uniform, although individual particles
keep moving. Diffusion is generally faster at higher temperatures.

\begin{center}% Replaces images/diffusion_diagram.png (original retained)
\resizebox{0.85\linewidth}{!}{\DiffusionWorksheetFigure}\end{center}

\begin{minipage}[c]{0.13\linewidth}
  \includegraphics[width=\linewidth]{images/image35.png}
\end{minipage}
\begin{minipage}[c]{0.83\linewidth}
  \href{https://phet.colorado.edu/sims/html/diffusion/latest/diffusion_all.html?locale=en}
  {\textcolor{DeepBlue}{\underline{Here}}}
  , an online simulation of particles at different
  temperatures can be found.
\end{minipage}\hfill


\textbf{Osmosis} is the net movement of a solvent through a
semipermeable membrane from a region of lower solute
concentration to a region of higher solute concentration.
The membrane allows the solvent to pass but restricts
the movement of certain dissolved particles.

\begin{center}% Replaces images/diffusion_diagram.png (original retained)
\resizebox{0.85\linewidth}{!}{\DiffusionWorksheetFigure}\end{center}


\begin{InfoBox}
\BoxHeading{Diffusion and Osmosis}
\textbf{Diffusion} is the net movement of particles from
higher to lower concentration.

\textbf{Osmosis} is the net movement of the solvent through
a semipermeable membrane towards the region of higher
solute concentration.
\end{InfoBox}


\Exercise[diff]{\textbf{Explain why diffusion is usually faster in gases than in liquids.}}

\ResponseGap

\HTwo{Mixtures}

Most materials encountered in everyday life are mixtures. They consist of two or more physically combined substances.


\CompareBox{\textbf{Pure substance}}{A pure substance has a fixed composition and characteristic properties. It is represented by one kind of particle. It cannot be separated into different substances by physical methods.}{\textbf{Mixture}}{A mixture contains at least two substances. Its composition can vary, and its components can be separated by physical methods because they retain different physical properties.}


\HThree{Homogeneous and Heterogeneous Mixtures}

A \textbf{homogeneous mixture} has a uniform composition
and appears as a single phase. Examples include air,
salt water, and brass.

A \textbf{heterogeneous mixture} contains distinguishable
regions or phases. Examples include soil, smoke, and milk.
Some heterogeneous mixtures, such as milk, may appear
uniform to the naked eye.

\par\smallskip

\begin{center}
  \scalebox{0.8}\HeterogeneousMixtureSketch
  \qquad
  \scalebox{0.8}\HomogeneousMixtureSketch
\end{center}

\HThree{Types of Mixtures}

\Exercise[mix-types]{Complete the table by writing the name of each mixture type and drawing a simple particle model.}



\begin{StandardTable}{
		colspec={X[1,c,m] X[2,c,m] X[1.15,l,m]},
		colsep=4pt,
		row{2-13}={ht=15mm}
	}
	
	{states of matter} & {name and particle model} & {example} \\
	
	& & rock \\
	& & brass, steel, bronze \\
	& & muddy water, ink, dispersion paint \\
	& & salt water, sugar water \\
	& & smoke from a pipe, dust-laden air in a sand storm \\
	& & sponge \\
	& & milk, salad dressing, hand cream \\
	& & Alcoholic drinks \\
	& & rain clouds (droplets of water in the air) \\
	& & soap bubbles \\
	& & Mineral water (contains gases like oxygen) \\
	& & air \\
	
\end{StandardTable}


\Exercise[mix-pure]{\textbf{Decide in which cases a pure substance is present.}}

\begin{ExerciseParts}
\item table salt solution
\item graphite
\item table salt
\item diamond
\item air
\item dry ice
\end{ExerciseParts}

\Exercise[mix-hethomo]{\textbf{Which of the following mixtures of substances are called heterogeneous and which are homogeneous? Also specify the matter of state of the components.}}

\begin{ExerciseParts}
\item Soil
\item Rose gold
\item Natural gas
\end{ExerciseParts}

\Needspace{45mm}
\Exercise{\textbf{Draw the following substances on a particle level}\textbf{.}}
\begin{center}% Replaces images/mixtures_ex20.png (original retained)
\resizebox{0.96\linewidth}{!}{\MixtureDrawingWorksheetFigure}\end{center}

\Needspace{95mm}
\Exercise[mix-assign]{\textbf{Name the state of matter for the individual particles in the mixtures in the figures below and name the mixtures.}}
\ExerciseOneTwentySevenMixturesFigure

\Exercise[irgend]{balbalbal}

\Exercise[mix-meri]{\textbf{What kind of mixture is a }meringue\textbf{ before and after baking? Does a chemical reaction take place during baking, or is it a change in the state of matter?}}

\ResponseGap



\HTwo{Separation Methods}

The components of a mixture can often be separated without a chemical reaction because each component retains its own physical properties. A suitable separation method \textbf{exploits a difference in properties} such as particle size, density, solubility, boiling point, or attraction to a surface.


\HThree{Separating a Mixture of Sand and Table Salt}

Sand and table salt can be separated in three steps. Each step is a physical process: no new substance is formed.


\HFour{1. Extraction}

Table salt dissolves easily in water:

\ResponseGap[170pt]


\HFour{2. Filtration}

\hspace{85pt} sub-microscopic level \hspace{110pt} macroscopic level

\begin{center}
% Replaces images/filtration_diagram.png (original retained)
\resizebox{0.85\linewidth}{!}{\FiltrationFigure}
\end{center}


\HFour{3. Evaporation}

\ResponseGap[100pt]

\HThree{Summary}

\begin{StandardTable}{
		colspec={
			X[1.7,l,m]
			X[2.3,l,m]
			X[1.3,l,m]
			X[1.1,l,m]
			X[1.5,l,m]
		},
		colsep=4pt,
		row{1}={ht=12mm}
	}

	{Separation method} &
	{Description} &
	{Properties used for separation} &
	{Type of mixture} &
	{Typical application} \\
	
	\textbf{distillation} &
	vaporization of substances with different boiling points
	followed by condensation &
	boiling point &
	l/l (hom.) &
	refinement of crude oil \\
	
	\textbf{Chromatography} &
	relies on the different ability of substances to adsorb
	to certain surfaces (stationary phase) and to dissolve
	in certain solvents (mobile phase) &
	solubility\newline
	particle size\newline
	adsorption\newline
	attraction &
	Several possibilities &
	Analysis and separation of mixtures in the lab \\
	
	\textbf{sieving} &
	separating particles of different sizes through
	the holes of a sieve &
	particle size &
	s/s &
	gravel quarry \\
	
	\textbf{filtration} &
	insoluble solids are retained by a filter through
	which only liquids can pass &
	particle size &
	s/l &
	coffee\newline tea \\
	
	\textbf{sedimentation} &
	particles in a suspension settle at the bottom
	of the fluid &
	density &
	s/l (het.) &
	sedimentation tanks in sewage treatment plants \\
	
	\textbf{decantation} &
	pouring off the liquid top layer &
	density &
	s/l\newline l/l (het.) &
	separation of mixtures of water and sand or oil and water \\
	
	\textbf{centrifugation} &
	making use of the centrifugal force &
	density &
	s/s\newline s/l\newline l/l &
	production of cream \\
	
	\textbf{extraction} &
	relies on the different solubilities of substances &
	solubility &
	s/s &
	production of salt\newline making of coffee or tea \\
	
	\textbf{evaporation} &
	evaporation of solvent &
	bp &
	s/l (hom.) &
	recovery of table salt
	
\end{StandardTable}